\metatitle{Symmetric functions as a vector space}
\metadescription{Vector-space and operator structures on symmetric functions,
including Fock space, the Hall inner product, the standard involution, adjoint
operators, and the Hopf algebra structure.}


\section[symmetricFunctionsSpace]{Vector-space properties of symmetric functions}

The ring of symmetric functions is also a graded vector space with several
natural extra structures.  This page collects the Fock-space viewpoint, the
Hall inner product, standard involutions and adjoint operators, and the Hopf
algebra structure on $\spaceSym$.


\subsection[fockSpace]{Fock space}

In the context of symmetric functions, \defin{Fock space} $\mathcal{F}$ is a
Hilbert space isomorphic to $\spaceSym$.
It provides a framework where bases of symmetric functions are viewed as states
created by operators acting on a vacuum.
This point of view is particularly powerful for proving identities involving
\hyperref[schurS]{Schur functions} via the \enquote{Boson--Fermion
correspondence}, and for connecting symmetric functions to integrable systems.
For background, see \cite[Chap.~14]{Kac1990} and
\cite[Chap.~4]{MiwaJimboDateReid2000}.

\name{Rapha{\"e}l Fesler}, \name{Marvin Anas Hahn},
\name{Maksim Karev}, and \name{Hannah Markwig} introduce a two-parameter
CJT-refinement of Jucys--Murphy theory acting on Fock space
\cite{FeslerHahnKarevMarkwig2025x}.  Their framework unifies Schur and zonal
actions, conjecturally includes Jack actions, and gives cut-and-join
recursions for $b$-Hurwitz numbers.

The space can be described in two equivalent ways:

\textbf{The Bosonic picture:} This picture identifies $\mathcal{F}$ with the
polynomial ring generated by the \hyperref[powerSum]{power-sum symmetric functions},
$\setC[\powerSum_1, \powerSum_2, \dotsc]$.
The \defin{Heisenberg algebra} acts on this space via operators $\alpha_n$ for
$n \in \setZ \setminus \{0\}$.
For $n \in \setZ \setminus \{0\}$, we define:
\[
\alpha_{-n} \cdot f = \powerSum_n \cdot f
\quad \text{(creation for $n > 0$)}, \qquad
\alpha_{n} \cdot f = n \frac{\partial f}{\partial \powerSum_n}
\quad \text{(annihilation for $n > 0$)}.
\]
These satisfy the commutation relation
$[\alpha_n, \alpha_m] = n \delta_{n+m,0} \mathrm{Id}$.
The \defin{vacuum state} $|0\rangle$ corresponds to the unit $1 \in \spaceSym$.
In this picture, the \hyperref[hallInnerProduct]{Hall inner product} arises
naturally, making the adjoint of multiplication by $\powerSum_n$ equal to
the differential operator $n \frac{\partial}{\partial \powerSum_n}$.


\textbf{The Fermionic picture:} Here, $\mathcal{F}$ is constructed as a
semi-infinite wedge space $\Lambda^{\frac{\infty}{2}} V$, where
$V = \bigoplus_{i \in \setZ} \setC v_i$.
A basis is given by semi-infinite wedges of the form
\[
  |S\rangle = v_{i_1} \wedge v_{i_2} \wedge v_{i_3} \wedge \dotsb
\]
where $i_1 > i_2 > \dotsb$ is a decreasing sequence of integers such that
$i_k = -k+1$ for $k$ large enough.
This sequence corresponds to a partition $\lambda$ via the relation
$\lambda_k = i_k - (-k+1)$.
Combinatorially, the set $\{i_1, i_2, \dotsc \}$ is often visualized as a
\defin{Maya diagram}.


\textbf{The Boson--Fermion correspondence:}
There is a canonical isomorphism
$\sigma: \mathcal{F}_{\text{Fermion}} \to \mathcal{F}_{\text{Boson}}$.
Under this map, the basis wedge corresponding to a partition $\lambda$ maps exactly
to the \hyperref[schurS]{Schur function} $\schurS_\lambda$:
\[
  \sigma(v_{i_1} \wedge v_{i_2} \wedge \dotsb ) = \schurS_\lambda(\xvec).
\]
This correspondence explains why Schur functions form an orthonormal basis:
they correspond to the orthonormal basis vectors of the exterior algebra.
Furthermore, the \hyperref[mapsOnSymmetricFunctions]{Bernstein creation operators}
can be viewed as the image of Fermionic creation operators acting on the Bosonic
space.


\subsection[symmetricGeneratingFamilies]{Generating families}

\name{Velmurugan S} gives criteria for when a homogeneous sequence
$u_n\in\Lambda_n$ is an algebraically independent set of generators for the
algebra of symmetric functions \cite{Velmurugan2024x}.  The results apply to
many standard families indexed by partitions or skew partitions, including
monomial, forgotten, skew complete, skew elementary, Schur, skew Schur,
Hall--Littlewood, $q$-Whittaker, Macdonald, and integral-form Macdonald
functions.

\subsection[hallInnerProduct]{Hall inner product}


Let $\langle \cdot, \cdot \rangle$ be the inner product on symmetric functions
such that $\langle \powerSum_\lambda, \powerSum_\mu \rangle = \delta_{\lambda\mu}z_\lambda$,
that is, the \hyperref[powerSum]{power-sum} functions form an orthogonal basis.
Then $\langle \schurS_\lambda, \schurS_\mu \rangle = \delta_{\lambda\mu}$,
so that the Schur polynomials form an orthonormal basis for $\spaceSym$.
See the \hyperref[prelimPartitionStatistics]{preliminaries} for the definition of $z_\lambda$.


The following properties hold for the Hall inner product:
\begin{itemize}

\item The \hyperref[monomial]{monomial basis} and the
\hyperref[completeH]{complete homogeneous basis} are \defin{dual}:
$\langle \monomial_\lambda, \completeH_\mu \rangle = \delta_{\lambda\mu}$.

\item Any symmetric function $f$ defines a map
$f:\spaceSym \to \spaceSym$ by multiplication.
This gives an \defin{adjoint map} $f^{\perp}: \spaceSym \to \spaceSym$
called \defin{skewing by $f$}.
That is,
\[
\langle f\cdot g, h \rangle = \langle g, f^{\perp} h \rangle.
\]

For example, multiplication by $\schurS_\mu$ is adjoint to the linear map that
sends $\schurS_\lambda$ to $\schurS_{\lambda/\mu}$.
In other words,
\[
\langle \schurS_\mu \schurS_\nu, \schurS_{\lambda} \rangle
=
\langle \schurS_\nu, \schurS_{\lambda/\mu} \rangle.
\]
We also have $\powerSum_n^{\perp} = n \frac{\partial}{\partial \powerSum_n}$.


\item The \hyperref[standardInvolution]{$\omega$ involution} is self-adjoint:
\[
\langle \omega(f), g \rangle = \langle f, \omega(g) \rangle
\]


\item Multiplication and comultiplication are adjoint; see
\hyperref[hopfAlgebraSymmetricFunctions]{the Hopf algebra structure} on
$\spaceSym$.

\item Schur positivity for an operator $T$ implies Schur positivity for the
adjoint.  More precisely,
\[
T(\schurS_\lambda) = \sum_{\mu} c_{\lambda \mu} \schurS_\mu
\iff
T^{\perp}(\schurS_\mu) = \sum_{\lambda} c_{\lambda \mu} \schurS_\lambda.
\]
This holds since the Schur functions constitute an orthonormal basis.
\end{itemize}

See also \hyperref[mapsOnSymmetricFunctions]{these other maps on symmetric functions}.


\subsection[standardInvolution]{Involution on symmetric functions}

There is a \defin{standard involution}\label{omegaInvolutionDefinition}
$\omega$ on symmetric functions, defined by
$\omega(\elementaryE_\lambda) = \completeH_\lambda$ for all $\lambda$.
It acts on the power-sum symmetric functions and the
\hyperref[schurS]{Schur polynomials} as
\[
\omega(\powerSum_\lambda)
=
(-1)^{|\lambda|-\length(\lambda)}\powerSum_\lambda
\quad \text{and} \quad
\omega(\schurS_\lambda) = \schurS_{\lambda'}.
\]

This involution has \hyperref[gesselProperties]{natural extensions} to
quasisymmetric functions by defining it on the
\hyperref[gessel]{Gessel quasisymmetric functions}.


\begin{lemma}[From \cite{BernardiNadeau2020}]
Let $f(t)$ be a formal power series (or polynomial) with $f(0)=1$.
Then
\[
\omega \left( f(x_1)f(x_2) \dotsm \right) = \frac{1}{f(-x_1)f(-x_2) \dotsm}.
\]
\end{lemma}


\subsection[mapsOnSymmetricFunctions]{Other maps on symmetric functions}


\begin{definition}[The Adams operator and the Verschiebung operator]

The \defin{$k$th Adams operator}\label{adamsOperatorSymmetricFunctions}
on $\spaceSym$ is the ring homomorphism
\[
 f(\xvec) \mapsto f[ \powerSum_k ] = f(x_1^k, x_2^k, \dotsc ),
\]
which is a \hyperref[plethysm]{plethysm}.

The \hyperref[hallInnerProduct]{adjoint operator} to the Adams operator
is the \defin{Verschiebung operator}\label{verschiebungOperatorSymmetricFunctions}
$\phi_k$, which is defined on the power-sum symmetric functions via
\[
 \phi_k \powerSum_m(\xvec) =
 \begin{cases}
k\cdot \powerSum_{m/k}(\xvec) & \text{ if } k \mid m \\
 0 & \text{ otherwise,}
 \end{cases}
\]
and then extended as a ring homomorphism.
\end{definition}


\begin{definition}[The Bernstein operator]\label{bernsteinOperator}
The \defin{Bernstein operator} $\mathbf{B}_r$ is defined as
\[
   \mathbf{B}_r \coloneqq \sum_{j \geq 0} (-1)^j \completeH_{r+j} \elementaryE^{\perp}_j.
\]
It is also called the \defin{creation operator} because it satisfies
\[
 \schurS_{\lambda} = \mathbf{B}_{\lambda_1} \mathbf{B}_{\lambda_{2}} \dotsb \mathbf{B}_{\lambda_{\ell}}(1).
\]
That is, we can create the Schur function $\schurS_{\lambda}$ from $1$, and
$\mathbf{B}_r$ is the operator that adds a new top row of length $r$.

\end{definition}


\begin{definition}[The hook filter]\label{def:hookFilter}
We define the linear map $\phi$ on the ring of quasisymmetric functions
by its action on the \hyperref[gesselDefinitionSubsets]{Gessel fundamental basis}.
For $S \subseteq [n]$,
\[
\phi(\gessel_{n+1,S}(\xvec)) =
\begin{cases}
t(t-1)^i & \text{if } S = \{i+1, i+2, \dotsc, n\}, \\
0 & \text{otherwise.}
\end{cases}
\]
\end{definition}

\begin{proposition}
For a partition $\lambda \vdash (n+1)$, the map $\phi$ acts as an exact filter
for hook shapes on \hyperref[schurS]{Schur functions}:
\[
\phi(\schurS_\lambda(\xvec)) =
\begin{cases}
t(t-1)^i & \text{if } \lambda = (i+1, 1^{n-i}), \\
0 & \text{otherwise.}
\end{cases}
\]
\end{proposition}

This map has an important role for the
\hyperref[chromaticQuasisymmetric]{chromatic symmetric functions}
\cite{Stanley1995}.
In particular, it has the property that
$\phi(\elementaryE_\mu) = t^{\length(\mu)}$.


\subsection[hopfAlgebraSymmetricFunctions]{Hopf algebra}

The algebra of symmetric functions admits a \hyperref[hopfAlgebra]{Hopf algebra}
structure,
where the coproduct is defined as
\[
\Delta( \elementaryE_k ) = \sum_{j=0}^k \elementaryE_j \otimes \elementaryE_{k-j}
\]
For introductions to Hopf algebras and symmetric functions, see
\cite{Cheng2014HopfSymmetricFunctions,GrinbergReiner2014Hopf}.


The Hopf algebra \defin{antipode}\label{symmetricFunctionAntipode} on symmetric functions
has the property that
\[
S( \schurS_{\lambda/\mu} ) = (-1)^{|\lambda/\mu|} \schurS_{\lambda'/\mu'}.
\]
For a combinatorial proof of this, see \cite{ChoHwangLee2026x}.
